\section*{Capítulo \ref{ch:g10:drainage-basins} --- Cuencas hidrográficas, ríos y crecidas}

\begin{solution}{exo:g10:drainage-basins:1}
La \omterm{def:g10:drainage-basins:drainage-basin}{cuenca hidrográfica} de un punto de un río es toda la superficie de tierra
cuya escorrentía va a parar al río aguas arriba de ese punto y pasa por él.
La \omterm{def:g10:drainage-basins:drainage-basin}{divisoria de aguas} es la línea, a lo largo de las crestas, que separa dos
cuencas vecinas. Dos gotas que caen a un metro una de otra, a ambos lados de
una divisoria, bajan por vertientes opuestas hacia dos cuencas distintas: sus
ríos pueden llegar a mares distintos.
\end{solution}

\begin{solution}{exo:g10:drainage-basins:2}
El \omterm{def:g10:drainage-basins:discharge}{caudal} es el volumen de agua que atraviesa una sección transversal del río
cada segundo, en \unit{m^3/s}. \qty{1}{m^3/s} son \qty{1000}{L} cada
segundo; en un día, $1000 \times 86\,400 = 86.4$ millones de litros.
\end{solution}

\begin{solution}{exo:g10:drainage-basins:3}
Antes de la crecida, el río llevaba unos \qty{330}{m^3/s} (entre unos
\num{240} y \qty{390}{m^3/s}). El pico fue de \qty{9231}{m^3/s} el
21 de enero. La subida empezó el 18 de enero (\qty{362}{m^3/s}): duró tres
días.
\end{solution}

\begin{solution}{exo:g10:drainage-basins:4}
$1/50 = 0.02$, un \qty{2}{\percent}. No: los años son independientes; la
probabilidad vuelve a ser del \qty{2}{\percent} este año, pase lo que pase el
año pasado.
\end{solution}

\begin{solution}{exo:g10:drainage-basins:5}
Un pie cúbico son $0.3048^3 \approx \qty{0.02832}{m^3}$, luego
$593\,000 \times 0.02832 \approx \qty{16800}{m^3/s}$. El Amazonas lleva
$206\,000/16\,800 \approx 12$ veces más.
\end{solution}

\begin{solution}{exo:g10:drainage-basins:6}
Áreas de las franjas: $1.5 \times 0.3 = 0.45$, $0.90$, $1.05$ y
\qty{0.60}{m^2}: $A = \qty{3.0}{m^2}$. \omterm{def:g10:drainage-basins:discharge}{Caudales}: $0.2 \times 0.45 = 0.09$, $0.5 \times 0.90
= 0.45$, $0.6 \times 1.05 = 0.63$, $0.3 \times 0.60 = 0.18$: $Q =
\qty{1.35}{m^3/s}$. Velocidad media: $1.35/3.0 = \qty{0.45}{m/s}$.
\end{solution}

\begin{solution}{exo:g10:drainage-basins:7}
$66 \times 10^6 \times 0.020 = \qty{1.32e6}{m^3}$. Dos días son
$2 \times 86\,400 = \qty{172800}{s}$: $1.32 \times 10^6/172\,800 \approx
\qty{7.6}{m^3/s}$ de media.
\end{solution}

\begin{solution}{exo:g10:drainage-basins:8}
Ordenados: 120, 88, 65, 59, 51, 47, 42, 38, 35. Con $n = 9$, $T = 10/m$: el
mayor (120) tiene $T = 10$ años; el segundo (88), 5 años; el tercero (65),
unos 3.3 años. La mediana, 51, tiene el rango 5 y $T = 2$ años: se iguala o
se supera un año de cada dos.
\end{solution}

\begin{solution}{exo:g10:drainage-basins:9}
El \omterm{def:g10:drainage-basins:hydrograph}{hidrograma} de la ciudad sube antes y mucho más alto, y baja deprisa: en
pendientes fuertes, tejados y calles, se infiltra poca agua y la escorrentía
llega al arroyo deprisa y toda a la vez; su \omterm{def:g10:drainage-basins:hydrograph}{tiempo de retardo} es corto. El
\omterm{def:g10:drainage-basins:hydrograph}{hidrograma} de la cuenca con bosque sube más tarde, hasta un pico más bajo y
más ancho: el suelo y el bosque absorben buena parte de la lluvia y la
sueltan despacio, y el agua corre lentamente por pendientes suaves. Su \omterm{def:g10:drainage-basins:hydrograph}{caudal
base} sigue más alto después de la tormenta.
\end{solution}

\begin{solution}{exo:g10:drainage-basins:10}
$9231 \times 86\,400 \approx \qty{7.98e8}{m^3}$, repartidos sobre
$29\,900 \times 10^6 = \qty{2.99e10}{m^2}$: unos \qty{0.027}{m},
\qty{27}{mm}. El \omterm{def:g10:drainage-basins:hydrograph}{caudal base} da $330 \times 86\,400/2.99 \times 10^{10}
\approx \qty{0.001}{m}$, alrededor de \qty{1}{mm} al día: el día del pico
pasó unas 28 veces más agua.
\end{solution}

\begin{solution}{exo:g10:drainage-basins:11}
$P = 1 - 0.9^{10} \approx 1 - 0.349 = 0.651$: alrededor del
\qty{65}{\percent}. No llega ninguna vez con una probabilidad de $0.9^{10} \approx 0.35$, alrededor del \qty{35}{\percent}.
\end{solution}

\begin{solution}{exo:g10:drainage-basins:12}
El \omterm{def:g10:drainage-basins:hydrograph}{hidrograma} es un triángulo de tres días de largo y \qty{1200}{m^3/s} de
alto: su área es $\tfrac12 \times 1200 \times 3 \times 86\,400 \approx
\qty{1.56e8}{m^3}$, \num{156} millones de metros cúbicos. La parte que pasa
de \qty{800}{m^3/s} es un triángulo más pequeño de la misma forma, tres veces
menor en altura ($400$ en vez de $1200$) y en duración (un día en vez de
tres): $\tfrac12 \times 400 \times 86\,400 \approx \qty{1.73e7}{m^3}$,
\num{17} millones de metros cúbicos. Es menos que los \num{60} millones que
puede almacenar el embalse: sí.
\end{solution}

\begin{solution}{exo:g10:drainage-basins:13}
Lleno hasta sus orillas, el río tiene $800/(100 \times 2) = \qty{4}{m}$ de
profundidad. (1) La llanura de inundación lleva \qty{400}{m^3/s} a lo largo
de \qty{900}{m} a \qty{0.5}{m/s}: su profundidad es de $400/(900 \times 0.5) \approx \qty{0.9}{m}$. (2) Entre los diques: $1200/(100 \times 2) = \qty{6}{m}$ de profundidad, \qty{2}{m} por encima de sus orillas, más del
doble de la profundidad en la llanura sin diques.
\end{solution}

\begin{solution}{exo:g10:drainage-basins:14}
Buscamos $(1 - 1/T)^{30} = 0.5$, luego $1 - 1/T = 0.5^{1/30} \approx 0.97716$,
$1/T \approx 0.02284$ y $T \approx 44$ años.
\end{solution}

\begin{solution}{exo:g10:drainage-basins:15}
Cada cuadrado manda su agua a su vecino más bajo (filas de norte a sur,
columnas de oeste a este):
\begin{center}
\begin{tabular}{|c|c|c|c|}
\hline
salida & oeste & oeste & oeste \\ \hline
norte & oeste & oeste & sur \\ \hline
norte & norte & sur & sur \\ \hline
norte & este & este & salida \\ \hline
\end{tabular}
\end{center}
Siguiendo las flechas, los cuadrados 47, 39, 44, 37, 46 y 30 desaguan en el
punto de salida del sureste: 6 cuadrados, \qty{0.06}{km^2}. Los otros 10
cuadrados desaguan en el punto de salida del noroeste: \qty{0.10}{km^2}. La
divisoria parte del borde este entre 50 y 47, y pasa luego entre 48 y 47,
entre 48 y 44, entre 49 y 44, entre 49 y 46, y entre 51 y 46 hasta el borde
sur.
\end{solution}

\begin{solution}{pb:g10:drainage-basins:1}
\textbf{1.} $11\,560 \times 1.609^2 \approx \qty{29900}{km^2}$.

\textbf{2.} La \omterm{def:g10:drainage-basins:drainage-basin}{divisoria de aguas} de la cuenca situada aguas arriba de la
estación.

\textbf{3.} $484\,000 \times 0.3048^3 \approx \qty{13700}{m^3/s}$.

\textbf{4.} $13\,705/3341 \approx 4.1$ veces.

\textbf{5.} $A = 400 \times 8 = \qty{3200}{m^2}$ y $v = Q/A = 13\,705/3200
\approx \qty{4.3}{m/s}$.

\textbf{6.} Unos \qty{330}{m^3/s}.

\textbf{7.} \qty{9231}{m^3/s}, el 21 de enero.

\textbf{8.} Tres días.

\textbf{9.} $9231 \times 86\,400/(29\,900 \times 10^6) \approx
\qty{0.027}{m}$: unos \qty{27}{mm}.

\textbf{10.} Llegó agua nueva al río: otra lluvia sobre la cuenca (o la
fusión de la nieve) mandó un segundo impulso de escorrentía, más pequeño.

\textbf{11.} Diecisiete días después del pico. La subida procede de la
escorrentía de las laderas cercanas a los arroyos, que llega deprisa y toda a
la vez; en la rama descendente, el río sigue recibiendo el agua de las partes
más lejanas de la cuenca y el agua que sale despacio del suelo.

\textbf{12.} $T = n + 1 = 96$ años.

\textbf{13.} $96/2 = 48$ años.

\textbf{14.} $96/3 = 32$ años.

\textbf{15.} $96/4 = 24$ años y $96/5 = 19.2$ años. No mucho: en ambos casos,
este \omterm{def:g10:drainage-basins:discharge}{caudal} vuelve más o menos una vez cada veinte años.

\textbf{16.} Cinco años (1936, 1937, 1943, 1972 y 1996): $5/95 \approx
0.053$, alrededor de un \qty{5}{\percent} al año, un \omterm{def:g10:drainage-basins:return-period}{periodo de retorno} de
unos $96/5 \approx 19$ años.

\textbf{17.} $1/100 = 0.01$, un \qty{1}{\percent}.

\textbf{18.} $1 - 0.01 = 0.99$.

\textbf{19.} $0.99^{30} \approx 0.740$.

\textbf{20.} $1 - 0.740 = 0.26$: \textbf{la probabilidad de al menos una
crecida de 100 años en 30 años es de alrededor del \qty{26}{\percent}}, una
entre cuatro.
\end{solution}
